%----------------------------------------------------------------------------------------
% OPMAAK VAN AXIOMATH: gedeeld door alle cursussen. Pas ze alleen hier aan.
% Compact handboek voor in de les. Uit dezelfde bron komen drie versies:
%   \def\versie{handboek}    kern, uitgewerkte voorbeelden, oefeningen zonder uitwerking
%   \def\versie{leerkracht}  handboek + didactische tips, antwoorden
%   \def\versie{zelfstudie}  handboek + extra uitleg, extra opgaven, antwoorden
%----------------------------------------------------------------------------------------
\providecommand{\versie}{handboek}
% Van de map van de cursus naar de hoofdmap van de repo (de cursus of maak-pdf.py zet dit al)
\providecommand{\hoofdmap}{../../../../}

\usepackage{etoolbox}
\newif\ifhandboek \newif\ifleerkracht \newif\ifzelfstudie \newif\ifuitwerking
\ifdefstring{\versie}{handboek}{\handboektrue}{}
\ifdefstring{\versie}{leerkracht}{\leerkrachttrue\uitwerkingtrue}{}
\ifdefstring{\versie}{zelfstudie}{\zelfstudietrue\uitwerkingtrue}{}

\usepackage[dutch]{babel}
\usepackage[a4paper,left=2.5cm,right=2.5cm,top=2.4cm,bottom=2.2cm,headsep=6mm,footskip=8mm,headheight=14pt]{geometry}

\usepackage{amsmath}
% Source Sans voor tekst en koppen, newtxsf voor schreefloze wiskunde (beide onder de SIL Open Font License)
\usepackage[T1]{fontenc}
\usepackage[default]{sourcesanspro}
\usepackage{newtxsf}
\renewcommand{\familydefault}{\sfdefault}
\linespread{1.1}
\usepackage{icomma}
\usepackage{ragged2e}

\usepackage{array,booktabs,colortbl}
\usepackage[inline]{enumitem}
\usepackage{multicol}
\usepackage{graphicx}
% \graphicspath breidt macro's niet uit: werk de lijst eerst uit
\edef\gptmp{{afbeeldingen/}{./}{\hoofdmap Ontwerp/omslagen/}{\hoofdmap Ontwerp/logo/}}
\expandafter\graphicspath\expandafter{\gptmp}
\usepackage{xcolor}
\usepackage{tikz}
\usetikzlibrary{arrows.meta,decorations.pathreplacing,calc}
\usepackage[most]{tcolorbox}
\usepackage{fancyhdr}
\usepackage{answers}
\usepackage{needspace}
% Zoals \needspace, maar zonder voorkeur (\penalty-100) om bij elke kop te breken: anders blijven pagina's half leeg
\makeatletter
\renewcommand{\needspace}[1]{\par\penalty0\begingroup\setlength{\dimen@}{#1}%
  \dimen@ii\pagegoal\advance\dimen@ii-\pagetotal
  \ifdim\dimen@>\dimen@ii\break\fi\endgroup}
\makeatother
\usepackage{refcount}

% Elk hoofdstuk heeft een kleur. Donker (hd) = uitleg, fel (hf) = opdrachten.
% Contrast nagekeken: wit op donker minstens 5,3 en inkt op fel minstens 7,0.
\definecolor{hd1}{HTML}{1257B0}\definecolor{hf1}{HTML}{5CB8FF}% blauw
\definecolor{hd2}{HTML}{0A7A5C}\definecolor{hf2}{HTML}{2EDDA5}% groen
\definecolor{hd3}{HTML}{B24A06}\definecolor{hf3}{HTML}{FFA62B}% oranje
\definecolor{hd4}{HTML}{B0216F}\definecolor{hf4}{HTML}{FF7DBD}% roze
\definecolor{hd5}{HTML}{6A35B8}\definecolor{hf5}{HTML}{BC96FF}% paars
\definecolor{hd6}{HTML}{BF2626}\definecolor{hf6}{HTML}{FF8A7A}% rood
\definecolor{hd7}{HTML}{0A6E8A}\definecolor{hf7}{HTML}{3FD6F0}% cyaan
\definecolor{hd8}{HTML}{6E6300}\definecolor{hf8}{HTML}{FFD93D}% olijf
\definecolor{inkt}{HTML}{12202B}
% Rubrieken, voor de labels van Koppel
\colorlet{rubriek-analyse}{hd1}\colorlet{rubriek-algebra}{hd5}\colorlet{rubriek-meetkunde}{hd3}
\colorlet{rubriek-discreet}{hd2}\colorlet{rubriek-data}{hd4}
% Kleur van de voorkaft en de oplossingen. Per graad een eigen kleur: voorkant-gegevens.tex (uit
% Ontwerp/voorkanten/maak-voorkanten.py) zet \omslagaccent en \omslagfel; zonder die macro's de 3de graad.
\providecommand{\omslagaccent}{0B6470}\providecommand{\omslagfel}{3FD6F0}
\definecolor{accent}{HTML}{\omslagaccent}
\definecolor{fel}{HTML}{\omslagfel}
\definecolor{grijs}{HTML}{5F6368}
\definecolor{lijngrijs}{HTML}{9AA0A6}
\newcommand{\kopletter}{\sffamily\bfseries}
\newif\ifopdrachtdeel % vanaf de Opdrachten: tab, kop en sectienummers in de felle tint
\DeclareEmphSequence{\color{accent}\bfseries,\itshape,\upshape}
\newcommand{\eur}{\texteuro{}}

%----------------------------------------------------------------------------------------
% ALINEA'S
%----------------------------------------------------------------------------------------
\setlength{\parindent}{0pt}
\setlength{\parskip}{5pt plus 1pt}
\AtBeginDocument{\RaggedRight\frenchspacing
  \thickmuskip=5mu plus 1mu\relax\medmuskip=4mu plus 1mu minus 2mu\relax
  \setlength{\abovedisplayskip}{5pt plus 1pt}\setlength{\belowdisplayskip}{5pt plus 1pt}%
  \setlength{\abovedisplayshortskip}{2pt}\setlength{\belowdisplayshortskip}{2pt}}
\tolerance=800
\hyphenpenalty=2000
\clubpenalty=10000 \widowpenalty=10000
\setlength{\mathindent}{3mm}
\raggedbottom
\setlength{\columnsep}{7mm}
\setlength{\multicolsep}{3pt}
\makeatletter
\def\@minipagerestore{\RaggedRight\setlength{\parskip}{3pt}}
\makeatother

\setlist[itemize,1]{label=\textcolor{accent}{\textbullet},leftmargin=4mm,itemsep=.1ex,topsep=.3ex,parsep=0pt}
\setlist[enumerate,1]{label=\textbf{\alph*.},leftmargin=6mm,labelsep=1.4mm,itemsep=.5ex,topsep=.6ex,parsep=0pt,beginpenalty=10000,midpenalty=10000}% deelvragen blijven samen
% Korte deelvragen naast elkaar, uitgelijnd (bv. rijtjes getallen). Lange deelvragen: enumerate.
\newenvironment{reeks}[1][3]{\begin{multicols}{#1}\begin{enumerate}[midpenalty=0]}{\end{enumerate}\end{multicols}}% midpenalty=0: de kolommen moeten kunnen breken

%----------------------------------------------------------------------------------------
% KOPPEN
%----------------------------------------------------------------------------------------
\makeatletter
\def\@makechapterhead#1{%
  \begin{tikzpicture}[remember picture,overlay]
    \fill[accent] (current page.north west) rectangle ($(current page.north east)+(0,-3.5cm)$);
    \node[anchor=west,inner sep=0pt,text=fel,font=\kopletter\fontsize{54}{54}\selectfont]
      at ($(current page.north west)+(2.5cm,-1.95cm)$) {\thechapter};
    \node[anchor=south west,inner sep=0pt,text=white,font=\sffamily\small]
      at ($(current page.north west)+(4.4cm,-1.35cm)$) {\cursustitel};
    \node[anchor=north west,inner sep=0pt,text=white,font=\kopletter\fontsize{26}{31}\selectfont]
      at ($(current page.north west)+(4.4cm,-1.5cm)$) {#1};
  \end{tikzpicture}%
  \par\vspace*{1.05cm}}
\def\@makeschapterhead#1{%
  {\parindent\z@\RaggedRight\normalfont\vspace*{-6mm}%
   {\kopletter\fontsize{21}{25}\selectfont #1\par}
   \vspace{1mm}{\color{accent}\hrule height 1.4pt}\vspace{3mm}}}
\renewcommand{\section}{\@startsection{section}{1}{\z@}
  {-4ex \@plus -1ex \@minus -.4ex}{1.2ex \@plus .2ex}
  {\normalfont\kopletter\large}}
\DeclareRobustCommand{\secnummer}[1]{%
  \ifopdrachtdeel\def\secvul{fel}\def\sectekst{inkt}\else\def\secvul{accent}\def\sectekst{white}\fi
  \tikz[baseline=(c.base)]\node[fill=\secvul,rounded corners=2pt,text=\sectekst,inner xsep=1.5mm,inner ysep=1.5mm,font=\sffamily\bfseries\large](c){#1};} 
% Start van de Opdrachten: nieuwe pagina met een felle band
\newcommand{\opdrachtenkop}{%
  \clearpage\global\opdrachtdeeltrue\thispagestyle{plain}%
  \begin{tikzpicture}[remember picture,overlay]
    \fill[fel] (current page.north west) rectangle ($(current page.north east)+(0,-2.6cm)$);
    \node[anchor=west,inner sep=0pt,text=inkt,font=\kopletter\fontsize{26}{31}\selectfont]
      at ($(current page.north west)+(2.5cm,-1.45cm)$) {Opdrachten};
  \end{tikzpicture}%
  \phantomsection\addcontentsline{toc}{opdracht}{Opdrachten}%
  \par\vspace*{2mm}}
% Felle balk voor Koppel en Test jezelf
\newcommand{\opdrachtsectie}[1]{%
  \par\addvspace{3.5ex}\needspace{10\baselineskip}\noindent
  \phantomsection\addcontentsline{toc}{opdracht}{#1}%
  \tikz\node[fill=fel,rounded corners=2.5pt,inner xsep=3mm,inner ysep=1.5mm,text width=\dimexpr\linewidth-6mm\relax,align=left,text=inkt,font=\kopletter\large]{#1};%
  \par\nobreak\addvspace{.5ex}}
% Sprongenfiguur in een theoriekader; het argument staat boven de sprongen, bv. +v
\newcommand{\sprongfiguur}[1]{%
  \begin{tikzpicture}[x=.7cm,y=1cm,>={Stealth[length=1.8mm]},font=\small]
    \draw[lijngrijs] (-.2,0) -- (9.9,0);
    \foreach \x/\t in {0/u_1,1.6/u_2,3.2/u_3,7.8/u_{n-1},9.4/u_n}{
      \fill[accent] (\x,0) circle (2pt);
      \node[below=1pt] at (\x,0) {$\t$};}
    \node at (5.5,0) [fill=accent!9!white,inner sep=1pt] {$\cdots$};
    \foreach \a/\b in {0/1.6,1.6/3.2,7.8/9.4}{
      \draw[->,accent,thick] (\a,.08) to[bend left=45] node[above=-1pt,text=black,font=\scriptsize] {$#1$} (\b,.08);}
    \draw[decorate,decoration={brace,amplitude=4pt},grijs] (0,.62) -- (9.4,.62)
      node[midway,above=3pt,text=black] {$n-1$ sprongen};
  \end{tikzpicture}}
% Nummers hangen in de linkermarge: bij een sectie in een kader, bij een subsectie zonder kader in de kleur van de kop
\renewcommand{\@seccntformat}[1]{%
  \ifstrequal{#1}{subsection}%
    {\llap{\csname the#1\endcsname\hspace{2mm}}}%
    {\llap{\secnummer{\csname the#1\endcsname}\hspace{2mm}}}}
\setcounter{secnumdepth}{2}
\renewcommand{\subsection}{\@startsection{subsection}{2}{\z@}
  {-3ex \@plus -.8ex \@minus -.3ex}{1ex \@plus .2ex}
  {\normalfont\kopletter\color{accent}\nobreak}}% \nobreak: \color maakt anders een breekpunt vlak na een sectiekop
\makeatother
\newcommand{\blokkop}[1]{\leavevmode{\kopletter\small\color{accent}#1}}
\pretocmd{\section}{
   \ifnum\value{section}>0 \newpage\fi %nieuwe section start op nieuwe pagina, behalve de eerste
   \needspace{9\baselineskip}}{}{} % kop en kader niet los van de rest
% Niet meteen na een sectiekop: \needspace geeft daar een breekpunt, en dan blijft de sectiekop alleen onderaan staan
\makeatletter
\pretocmd{\subsection}{\if@nobreak\else\needspace{7\baselineskip}\fi}{}{}
\makeatother

\pagestyle{fancy}
\renewcommand{\chaptermark}[1]{\markboth{\thechapter.\ #1}{}}
\renewcommand{\sectionmark}[1]{}
\fancyhf{}
\fancyhead[L]{\sffamily\footnotesize\color{grijs}\leftmark\ifopdrachtdeel\ \textperiodcentered\ Opdrachten\fi}% rustig, zoals het paginanummer
\fancyfoot[R]{\sffamily\footnotesize\color{grijs}\thepage}
\renewcommand{\headrulewidth}{0pt}
\fancypagestyle{plain}{\fancyhf{}\fancyfoot[R]{\sffamily\footnotesize\color{grijs}\thepage}\renewcommand{\headrulewidth}{0pt}}

% Tab aan de rand van elke pagina: kleur, nummer en plaats tonen het hoofdstuk.
% 14 mm breed, zodat er genoeg overblijft als een printer de rand niet bedrukt.
\newlength{\tabx}\newlength{\taby}
\AddToHook{shipout/background}{%
  \ifnum\value{chapter}>0
    \setlength{\tabx}{\paperwidth}\addtolength{\tabx}{-14mm}%
    \setlength{\taby}{\dimexpr -4.3cm - 3cm*\numexpr\value{chapter}-1\relax\relax}%
    \ifopdrachtdeel\def\tabvul{fel}\def\tabtekst{inkt}\else\def\tabvul{accent}\def\tabtekst{white}\fi
    \put(\tabx,\taby){\textcolor{\tabvul}{\rule[-2.7cm]{14mm}{2.7cm}}}%
    \put(\tabx,\taby){\raisebox{-1.6cm}[0pt][0pt]{\makebox[8.5mm]{\color{\tabtekst}\kopletter\LARGE\thechapter}}}%
  \fi}

%----------------------------------------------------------------------------------------
% BLOKKEN
%----------------------------------------------------------------------------------------
\tcbset{blok/.style={enhanced,sharp corners,boxrule=0pt,frame hidden,
    left=4mm,right=3mm,top=2mm,bottom=2mm,before skip=2ex,after skip=2ex}}

% Theorie: definities, eigenschappen, formules en hun bewijs. Volledig omkaderd en getint.
\newtcolorbox{theorie}{enhanced,breakable,sharp corners,boxrule=1pt,colframe=accent,colback=accent!9!white,
  left=4mm,right=3mm,top=2mm,bottom=2mm,before skip=2ex,after skip=2ex,
  before upper={\RaggedRight\setlength{\parskip}{3pt}}}
% Een kader begint niet met een paar regels onderaan een pagina: het mag alleen verderop breken, bij een lang bewijs
\AtBeginEnvironment{theorie}{\needspace{10\baselineskip}}
% Formule in een theoriekader: vet en in de kleur van het hoofdstuk. Gebruik in wiskundemodus: $\formule{...}$
\newcommand{\formule}[1]{\text{\boldmath\color{accent}$\displaystyle #1$}}
% Bewijs: staat in het theoriekader, onder de eigenschap. Binnenin meestal de omgeving stappen.
% Binnenin: korte bindtekst met de formules op een eigen regel (\bewijsregel). Achteraan een vierkantje in de kleur van het hoofdstuk.
% Een bewijs is één blok (vbox): een kader breekt dus nooit midden in een bewijs, alleen ervoor.
\newenvironment{bewijs}{\par\addvspace{1.5mm}\vbox\bgroup\blokkop{Bewijs}\par\nobreak}{\par\nobreak\vspace{.6mm}\noindent\hfill\textcolor{accent}{\rule{2.2mm}{2.2mm}}\par\egroup}
\newcommand{\bewijsregel}[1]{\par\vspace{.6ex}\hspace*{6mm}$\displaystyle #1$\par\vspace{.6ex}}
% Stelling of eigenschap: eigen wit kader binnen het theoriekader, met de naam op de bovenlijn.
% De naam staat met haar midden op de lijn, en de lijn is daar onderbroken. Gebruik: \begin{stelling}{Naam} ... \end{stelling}
\newtcolorbox{stelling}[1]{enhanced,sharp corners,colback=white,colframe=accent,boxrule=.8pt,
  left=3mm,right=3mm,top=3.2mm,bottom=2mm,after skip=1.5ex,
  % geen breekpunt tussen een definitie en de stelling eronder; bovenaan het kader (prevdepth -1000pt) geen extra ruimte
  before={\par\ifdim\prevdepth=-1000pt\else\nobreak\addvspace{3.2ex}\fi\noindent},
  before upper={\RaggedRight\setlength{\parskip}{3pt}},
  title={#1},coltitle=accent,fonttitle=\kopletter\small,
  attach boxed title to top left={xshift=4mm,yshift=-\tcboxedtitleheight/2},
  boxed title style={enhanced,frame hidden,interior hidden,sharp corners,boxrule=0pt,
    left=1.5mm,right=1.5mm,top=.4mm,bottom=.4mm},
  underlay boxed title={%
    \fill[accent!9!white] (title.north west) rectangle (title.east);
    \fill[white] (title.west) rectangle (title.south east);}}
% Werkwijze: hoe je een soort opgave aanpakt. Enkel een balk links, geen tint, met een klein label.
\newtcolorbox{werkwijze}{blok,colback=white,
  borderline west={2.5pt}{0pt}{accent},
  before upper={\RaggedRight\setlength{\parskip}{3pt}\blokkop{Werkwijze}\quad}}

% Voorbeeld: altijd uitgewerkt, ook in het handboek. Oefeningen zijn dat niet.
\newcounter{voorbeeld}[chapter]
\providecommand{\theHvoorbeeld}{\thechapter.\arabic{voorbeeld}}
\newcommand{\voorbeeld}[3]{
  \par\addvspace{2ex}\refstepcounter{voorbeeld}%
  \vbox{%
    \setbox0=\hbox{%
      \tikz[baseline=(c.base)]\node[fill=accent,rounded corners=2pt,text=white,
        inner xsep=1.6mm,inner ysep=.6mm,
        font=\sffamily\bfseries\footnotesize](c)
        {Voorbeeld \thevoorbeeld\ifstrempty{#1}{}{: #1}};\quad}%
    \hangindent=\wd0 \hangafter=1
    \noindent\box0 #2\par
    \uitwerkingblok{#3}}%
  \addvspace{1.5ex}}

\newcommand{\uitwerkingblok}[1]{%
  \begin{tcolorbox}[blok,colback=white,borderline west={1pt}{0pt}{accent},before skip=.5ex,after skip=1ex,
      before upper=\RaggedRight\setlength{\parskip}{3pt}]
    #1
  \end{tcolorbox}}

% Uitleg bij een voorbeeld: cursief, net als de uitleg naast de stappen.
\newcommand{\uitleg}[1]{{\itshape #1\par}\vspace{.8ex}}

\newenvironment{stappen}{%
  \par\renewcommand{\arraystretch}{1.25}%
  \noindent\begin{tabular}{@{}>{$\displaystyle}p{7cm}<{$}@{\hspace{4mm}}>{\RaggedRight\small\itshape\arraybackslash}p{7.6cm}@{}}%
}{\end{tabular}\par}

% Alleen voor de leerkracht
\newcommand{\lkr}[1]{\ifleerkracht
  \begin{tcolorbox}[blok,colback=black!5!white,borderline west={1pt}{0pt}{grijs},before skip=.6ex,after skip=.8ex,
      before upper=\RaggedRight\sffamily\footnotesize]
    {\bfseries Voor de leerkracht.} #1
  \end{tcolorbox}\fi}

% Alleen in de zelfstudieversie: uitleg in lopende tekst
\newcommand{\thuis}[1]{\ifzelfstudie\par #1\par\fi}

% Alleen in de zelfstudieversie: extra opgave, met het antwoord onderaan de pagina
\newcounter{nujij}[chapter]
\providecommand{\theHnujij}{\thechapter.\arabic{nujij}}
\makeatletter
\newcommand{\nujij}[2]{\ifzelfstudie
  \refstepcounter{nujij}%
  \begin{tcolorbox}[blok,colback=white,frame hidden,borderline={1.2pt}{0pt}{fel},left=2.5mm,right=2.5mm,
      before upper=\RaggedRight]
    \blokkop{Nu jij \thenujij}\quad #1
  \end{tcolorbox}%
  \begingroup\long\def\@makefntext##1{\noindent\footnotesize ##1}%
  \footnotetext{\RaggedRight\textbf{Nu jij \thenujij:} #2}\endgroup\fi}
\makeatother

\newcommand{\vak}{\mbox{\fboxsep=0pt\fboxrule=.5pt\fbox{\rule{0pt}{1.8ex}\hspace{1.4em}}}}

%----------------------------------------------------------------------------------------
% OEFENINGEN: een kolom, nummer links, balkjes voor de moeilijkheid (bij de Opdrachten)
%----------------------------------------------------------------------------------------
\newcounter{oefening}[chapter] % nummering herstart per hoofdstuk
\providecommand{\theHoefening}{\thechapter.\arabic{oefening}}
\newcommand{\oefstrook}[1]{%
  \needspace{10\baselineskip}\par\noindent
  \tikz[baseline=(s.base)]\node[fill=fel,rounded corners=2.5pt,inner xsep=2.2mm,inner ysep=1mm,text=inkt,font=\sffamily\bfseries\small](s){#1};%
  \hspace{1.5mm}\textcolor{fel}{\leaders\hrule height 2.9pt depth -1.2pt\hfill\kern0pt}\par\nobreak}
\newcommand{\oefnummer}[1]{%
  \tikz[baseline=(c.base)]\node[fill=fel,rounded corners=2pt,inner xsep=1.2mm,inner ysep=.95mm,minimum width=7mm,text=inkt,font=\sffamily\bfseries\small](c){#1};}
\newenvironment{oefeningen}[1][]{\par\addvspace{1.2ex}\ifstrempty{#1}{}{\addvspace{1ex}\oefstrook{#1}}}{\par\addvspace{1ex}}
\newenvironment{oefening}[1][1]{%
  \needspace{5\baselineskip}% opgave niet los van de deelvragen
  \refstepcounter{oefening}%
  \renewcommand{\Currentlabel}{\theoefening}%
  \begin{list}{}{\setlength{\leftmargin}{15mm}\setlength{\labelwidth}{15mm}\setlength{\labelsep}{0pt}%
      \setlength{\topsep}{1.7ex}\setlength{\partopsep}{0pt}\setlength{\parsep}{3pt}\setlength{\itemsep}{0pt}}%
  \interlinepenalty=10000 % een oefening breekt niet over twee pagina's
  \item[{\makebox[15mm][l]{\oefnummer{\theoefening}\ifopdrachtdeel\hspace{.9mm}\niveau{#1}\fi}}]\ignorespaces
}{\end{list}}

% Kleine prent rechts naast een opdracht: \begin{metprent}{naam} opgave en deelvragen \end{metprent}
% De prent staat in afbeeldingen/prenten/, de solution komt na \end{metprent}. Optioneel: [breedte].
\newenvironment{metprent}[2][2.1cm]{%
  \def\prentnaam{#2}\def\prentbreedte{#1}%
  \begin{minipage}[t]{\dimexpr\linewidth-#1-4mm\relax}%
}{%
  \end{minipage}\hfill
  \begin{minipage}[t]{\prentbreedte}
    \vspace{-3mm}\raggedleft
    \includegraphics[width=\linewidth,height=2.1cm,keepaspectratio]{prenten/\prentnaam}
  \end{minipage}\par}

% Moeilijkheid: drie oplopende balkjes, niveau 1, 2 of 3. Ook duidelijk in zwart-wit.
\newcommand{\niveau}[1]{%
  \begin{tikzpicture}[baseline=.1mm,x=1mm,y=1mm]
    \foreach \b/\h in {1/1.6,2/2.8,3/4}{
      \ifnum\b>#1\relax\def\balk{black!20}\else\def\balk{inkt}\fi
      \fill[\balk] ({(\b-1)*1.4},0) rectangle ++(1.05,\h);}
  \end{tikzpicture}}
% Onderwerp van een opdracht, enkel in de leerkrachtenversie
\newcommand{\thema}[1]{\ifleerkracht{\sffamily\scriptsize\color{grijs}[#1]}\ \fi\ignorespaces}
% Label naar andere leerstof: de kleur toont de rubriek (analyse, algebra, meetkunde, discreet, data)
\newcommand{\koppel}[2]{%
  \tikz[baseline=(k.base)]\node[fill=rubriek-#1,rounded corners=2mm,text=white,inner xsep=2mm,inner ysep=.7mm,font=\sffamily\bfseries\footnotesize](k){%
    \tikz[baseline=-.2mm,x=1mm,y=1mm,line width=.5pt,white]{\draw[rounded corners=.5mm] (0,0) rectangle (1.9,1.1);\draw[rounded corners=.5mm] (1.3,0) rectangle (3.2,1.1);}\ #2};\ \ignorespaces}
\newcommand{\voorbeeldpil}[1]{%
  \tikz[baseline=(c.base)]\node[fill=accent,rounded corners=2pt,text=white,inner xsep=1.6mm,inner ysep=.6mm,font=\sffamily\bfseries\footnotesize](c){#1};}
\newcommand{\oefstrookje}[1]{%
  \tikz[baseline=(s.base)]\node[fill=fel,rounded corners=2.5pt,inner xsep=2.2mm,inner ysep=1mm,text=inkt,font=\sffamily\bfseries\small](s){#1};}

% Schema 'Fout? Herhaal'. In het hoofdstuk: \herhaalschema{\fv{test}{voorbeeld}{oefening} \fzv{test}{oefening} ...}
% met labels als argument. Het komt ook in de lijsten met oplossingen, met de nummers al ingevuld.
\protected\def\foutrij#1#2#3{\oefnummer{#1} & $\rightarrow$ & \ifstrempty{#2}{}{\voorbeeldpil{Voorbeeld #2}\ }\oefnummer{#3}\\}
\protected\def\foutschema#1{%
  \par\addvspace{2ex}\needspace{8\baselineskip}\noindent\blokkop{Fout? Herhaal}\par\nobreak\vspace{.6ex}%
  {\renewcommand{\arraystretch}{1.5}\noindent\begin{tabular}{@{}l@{\hspace{2.5mm}}c@{\hspace{2.5mm}}l@{}}#1\end{tabular}\par}\addvspace{1.5ex}}
\newcommand{\herhaalschema}[1]{%
  \begingroup
  \def\fv##1##2##3{\foutrij{\getrefnumber{##1}}{\getrefnumber{##2}}{\getrefnumber{##3}}}%
  \def\fzv##1##2{\foutrij{\getrefnumber{##1}}{}{\getrefnumber{##2}}}%
  \edef\schemaregels{#1}%
  \foutschema{\schemaregels}%
  \Writetofile{oplossingen}{\foutschema{\schemaregels}}%
  \Writetofile{uitgewerkteoplossingen}{\foutschema{\schemaregels}}%
  \endgroup}

%----------------------------------------------------------------------------------------
% OPLOSSINGEN
%----------------------------------------------------------------------------------------
\Newassociation{solution}{Sol}{oplossingen}
\Newassociation{fullsolution}{Fullsol}{uitgewerkteoplossingen}
\renewenvironment{Sol}[1]{\par\noindent\textcolor{accent}{\textbf{#1.}}~\ignorespaces}{\par\vspace{.4ex}}
\renewenvironment{Fullsol}[1]{\par\medskip\noindent\textcolor{accent}{\textbf{Oefening #1}}\par\nobreak}{\par}
\newcommand{\oplossingenkop}[2]{%
  \colorlet{accent}{hd#1}\colorlet{fel}{hf#1}%
  \par\vspace{2mm}\noindent{\kopletter\color{accent}#1.\ \ #2}\par\nobreak\vspace{1mm}}
\newcommand{\mychapter}[2]{%
  \clearpage % eerst de vorige pagina afsluiten: die houdt de kleur van het vorige hoofdstuk
  \global\opdrachtdeelfalse
  \colorlet{accent}{hd\number\numexpr\value{chapter}+1\relax}%
  \colorlet{fel}{hf\number\numexpr\value{chapter}+1\relax}%
  \Writetofile{oplossingen}{\protect\oplossingenkop{#1}{#2}}%
  \Writetofile{uitgewerkteoplossingen}{\protect\oplossingenkop{#1}{#2}}%
  \chapter{#2}}

%----------------------------------------------------------------------------------------
% VOORKAFT: in een eigen bestand, zodat een wijziging de bewaarde grafieken niet ongeldig maakt
%----------------------------------------------------------------------------------------
\input{\hoofdmap Ontwerp/voorkanten/voorkant}

%----------------------------------------------------------------------------------------
% INHOUD: hoofdstukken in hun eigen kleur, secties, en per hoofdstuk Opdrachten, Koppel en Test
%----------------------------------------------------------------------------------------
\makeatletter
\setcounter{tocdepth}{1} % onderdelen zonder nummer niet: de inhoud past zo op één pagina
\newcommand{\tochoofdstuk}{1}
\newcommand{\tocpuntjes}{\leaders\hbox to 2.4mm{\hss\textcolor{lijngrijs}{.}\hss}\hfill}
\def\toclevel@opdracht{1}
\def\l@chapter#1#2{\par\addvspace{4mm}%
  {\noindent\kopletter\large\color{inkt}%
   \def\numberline##1{\gdef\tochoofdstuk{##1}%
     \makebox[10mm][l]{\tikz[baseline=(c.base)]\node[fill=hd##1,rounded corners=2pt,text=white,inner xsep=1.6mm,inner ysep=.7mm](c){##1};}%
     \color{hd##1}}%
   #1\hfill #2\par}\nobreak\vspace{.6mm}}
\def\l@section#1#2{%
  {\noindent\hspace*{10mm}%
   \def\numberline##1{\makebox[9mm][l]{\textcolor{hd\tochoofdstuk}{\bfseries\small ##1}}}%
   #1\ \tocpuntjes\ #2\par}}
\def\l@opdracht#1#2{%
  {\noindent\hspace*{10mm}\makebox[9mm][l]{\textcolor{hf\tochoofdstuk}{\rule{4.5mm}{2.2mm}}}%
   {\small #1}\ \tocpuntjes\ {\small #2}\par}}
\newcommand{\inhoudstafel}{%
  \chapter*{Inhoud}%
  \begingroup\setlength{\parskip}{.8pt}\@starttoc{toc}\endgroup
  \clearpage}
\makeatother

\usepackage[hyperfootnotes=false]{hyperref}
\hypersetup{hidelinks,pdftitle={\cursustitel},pdfauthor={\cursusauteur},pdflang={nl}}
